\chapter{Primes in short intervals}\label{ch:oc-short}

This chapter asks how many primes lie in a short interval, from the Hardy--Littlewood conjectures to Gallagher's Poisson law and the Montgomery--Soundararajan variance.

\section{The Hardy--Littlewood conjectures}

The Hardy--Littlewood prime-tuples conjecture predicts the asymptotic count of prime $k$-tuples. Short-interval statistics are the same conjecture integrated over the starting point.

\section{Gallagher's Poisson law}

Gallagher showed that, on the Hardy--Littlewood conjectures, the number of primes in $(x,x+\lambda\log x]$ converges to a Poisson law of mean $\lambda$ as $x\to\infty$.

\section{The Montgomery--Soundararajan variance}

Montgomery and Soundararajan predicted that the variance of the number of primes in short intervals is smaller than the Poisson variance, by a factor about $\log(y/\log x)$ when the length $y$ grows. The aim is to see this as a repulsion between primes, not as independence.
